If a two-coordinate table is required and both stator axes may be shifted, F is the most convincing simple k-based compromise in this comparison. It retains two constant parameters and a globally simple linear coordinate map, improves joint interior prediction over A in every dataset, and is close to the best restricted nonlinear candidates. Its joint RMSE reductions relative to A are approximately 8.5\%, 10.4\% and 6.3\%. These percentages compare the defined flux-error statistic; they are not the derivative-angle measure. No application acceptance threshold is inferred.

E deserves consideration when the q-axis must remain the physical q-current and performance on the first machine configuration is the priority. It has the lowest reduced-model error in Set 1 and nearly matches the five-parameter G in Set 2. That result does not transfer unchanged to Set 3, where the simple two-axis F performs better. G has the lowest reduced-model joint error in Sets 2 and 3, but its extra coefficients, nonlinear inverse conditions and component trade-offs buy only a small improvement over the simpler nearest competitors. A remains a useful named baseline, while C is a better-founded constant d-shift fitting criterion rather than a universally superior predictor.

When flux accuracy dominates storage, H is preferable within the supported measured domain. It reduces the joint missing-group RMSE to approximately 1.5--1.8 mVs compared with approximately 5.4--7.0 mVs for the reduced candidates. This advantage costs an explicit excitation axis: the realized full tables have 137,781 nodes versus 6,561 for the reduced tables, or about 2.20 MB versus 0.105 MB for two double-precision flux arrays before masks and axes. It is an empirical reference with finite coverage, not an exact physical truth or an extrapolation guarantee.

The mathematical conclusion is unchanged by estimator improvements: a common constant d-shift requires both equations~\eqref{eq:both}, with the stated global fibre conditions. A variable pure d-shift cannot compensate an absent common local direction merely by adding scalar parameters. The independent-surface and derivative checks address that compatibility from different numerical paths, while reconstruction and whole-group prediction quantify its practical consequences. Temperature, sensor and preprocessing uncertainty prevent a unique physical error partition. The recommended reduced model is therefore an approximation to be assessed against an application's own flux-error requirements, not a lossless replacement established by these experiments.
